\section{Root differentials and finite covers}\label{in:foundations}

We allow, throughout the preparatory arguments, that the reduced divisor
\(D\) is empty or is a disjoint union of smooth divisors.  Its nonempty
components all have order two.  Write \(\mathcal X\) for the resulting
root orbifold and \(\pi:\mathcal X\to X\) for its coarse map.  At a point
of \(D\), a chart has coordinates
\[
 z_1=w_1^2,\qquad z_j=w_j\quad(j>1),
\]
with the sign action on \(w_1\).  For a prime divisor \(E\) of \(X\), set
\(r_E=2\) when \(E\subset D\), and \(r_E=1\) otherwise.  Saturation of a
meromorphic line direction in a locally free sheaf is understood on
these charts, equivariantly.

\begin{lemma}[Auxiliary projective source]\label{in:auxiliary}
Let \(X\) be a connected smooth projective variety with the root datum
just described.  There is a connected smooth projective variety \(T\)
and a surjective generically finite morphism \(d:T\to X\), Galois on
function fields and equivariant for its finite Galois group, such that
\(d^*D\) is twice an integral divisor.  This includes every exceptional
prime of \(d\).  The morphism lifts to \(T\to\mathcal X\), and induces a
homomorphism
\[
 \chi_d:\pi_1(T)\longrightarrow \pi_1^{\mathrm{orb}}(\mathcal X)
\]
with finite-index image, compatible with the map on divisor
complements.  A deck transformation changes this homomorphism only by
the base-point conjugation coming from the target orbifold group.
\end{lemma}

\begin{proof}
Choose a line bundle \(L\) such that \(L^{\otimes2}(-D)\) is very ample,
and choose finitely many sections \(b_i\) of this bundle having no
common zero on \(D\).  If \(s_D\) is the canonical section of
\(\mathcal O_X(D)\), adjoin square roots of the sections \(s_Db_i\) of
\(L^{\otimes2}\), interpreted after a rational trivialization of \(L\).
Their compositum is a finite Galois field extension.  Normalize \(X\)
in this field and take a projective resolution equivariant for the
finite group, unchanged over the smooth locus; denote the composite
by \(d\).  Canonical characteristic-zero resolution is functorial
for smooth morphisms and equivariant for group actions
\cite[Theorems~1.0.2--1.0.3]{Wlodarczyk}.  For projectivity, embed the
normalization equivariantly in projective space by taking the product
of the translates of a projective embedding and then the Segre
embedding.  The embedded resolution consists of blowups in invariant
smooth centres; its ambient space and final strict transform are
projective.  The normalization is integral, so its resolution is
connected.

Let \(F\) be any prime of \(T\) whose image is contained in \(D\).  At
the generic point of its image some \(b_i\) is a unit.  In local frames,
the square-root equation gives
\[
 2\operatorname{ord}_F(t_i)
   =\operatorname{ord}_F(d^*s_D).
\]
The roots \(t_i\) are regular on the normalization, since they are
integral over its local rings, and remain regular after resolution.
The equality proves evenness also when \(F\) is exceptional.  As the
components of \(D\) are disjoint, each \(d^*D_j\) is separately twice
an integral effective divisor.  The line bundle and canonical section
of its half divisor are a square root of the pulled-back divisor
bundle and section.  These data give the asserted map to the root
orbifold.  For empty \(D\) one may take the identity.

On \(T\setminus d^{-1}D\) the map to \(X\setminus D\) induces a map
to the orbifold group.  Every meridian filled in on passing to \(T\)
maps either trivially or to an even power of a meridian about one
smooth branch of \(D\).  The normal-meridian description proved in
Lemma~\ref{in:topology} therefore gives \(\chi_d\).  On suitable dense
opens \(d\) is a finite \'{e}tale cover.  Its group image has finite
index in the complement group, hence also in its orbifold quotient.
This image is contained in \(\operatorname{im}\chi_d\).  Finally,
transport between two sheets projects to a loop in \(X\setminus D\);
compatibility there gives the stated conjugation formula, and the
surjection from the open group to \(\pi_1(T)\) extends it to the
compact source.
\end{proof}

\begin{lemma}[Saturated lines and the differential bound]
\label{in:root-bound}
A meromorphic line direction in \(\Omega^p_{\mathcal X}\) has a
saturated holomorphic line bundle with a nonzero map to
\(\Omega^p_{\mathcal X}\).  If a holomorphic line bundle
\(\mathcal L\) has such a map, then
\(\kappa(\mathcal X,\mathcal L)\le p\).  Saturating the image does not
decrease its Iitaka dimension.  In particular, specialness of
\(\mathcal X\) implies specialness of \(X\).
\end{lemma}

\begin{proof}
Intersect the meromorphic line with the locally free differential
sheaf.  This saturated rank-one subsheaf is reflexive: it is the
kernel of a morphism from a locally free sheaf to a torsion-free
sheaf, and the depth criterion gives the reflexive condition on a
smooth chart.  A rank-one reflexive sheaf on a smooth chart is
invertible.  Uniqueness of saturation supplies equivariant gluing.
Replacing a line map by its saturation permits an effective
divisorial enlargement and therefore cannot decrease the section-map
dimension.

Pull \(\mathcal L\) and its map to the source \(T\) of
Lemma~\ref{in:auxiliary}.  The induced map to \(\Omega_T^p\) is
nonzero, since \(d\) is generically \'{e}tale.  Ratios of sections keep
their transcendence degree under this dominant map.  The manifold
differential bound in Assumption~\ref{in:input-special}, applied to
the smooth projective \(T\), consequently implies
\(\kappa(\mathcal X,\mathcal L)\le p\).  For the last assertion,
ordinary differential line maps pull back to root charts, and
\(H^0(\mathcal X,\pi^*M^{\otimes j})=H^0(X,M^{\otimes j})\) for an
ordinary line bundle \(M\).  An ordinary witnessing line would
therefore witness nonspecialness on \(\mathcal X\).
\end{proof}

We make precise the specialness convention used when applying the
orbifold literature.  For a simple normal crossing root datum of
finite orders \(m_i\), write locally \(z_i=w_i^{m_i}\).  A monomial
in \(\operatorname{Sym}^j\Omega^p_X\) in which \(dz_i\) occurs
\(t_i\) times has an invariant holomorphic lift exactly when its
coefficient has pole order at most
\begin{equation}\label{in:rounding}
 \left\lfloor t_i(1-m_i^{-1})\right\rfloor
 \quad\text{along }z_i=0.
\end{equation}
Indeed a coefficient of order \(-a_i\) pulls back with exponent
\(-m_i a_i+(m_i-1)t_i\) in \(w_i\).  These are precisely the local
generators of Campana's orbifold symmetric differential sheaves.

For a coarse generic differential line, saturate upstairs to obtain
\(\mathcal L\).  At every chart prime its generator is primitive,
as is each tensor power.  Thus the invariant sections of
\(\mathcal L^{\otimes j}\) are exactly the saturation of the
corresponding coarse generic power in the sheaf given by
\eqref{in:rounding}; equality across codimension two follows by
reflexivity.  Conversely an orbifold differential line determines
such a coarse generic line.  The invariant sheaves are coherent and
algebraic on the projective coarse space \cite{GAGA}.  A power
divisible by all stabilizer orders has trivial stabilizer character
and descends to an ordinary line bundle: locally average a frame
whose value at the point is invariant.  Such a frame remains
nonvanishing after shrinking.  Divisible powers compute the same
Iitaka dimension, since ratios can be raised to common powers.
The usual section-growth definition agrees as well; any nonzero
residue degree can be shifted into divisible degrees by multiplying
by a fixed nonzero section in the complementary residue degree,
obtained as a power of an existing section.  This identifies the
line-bundle definition used here with the saturated Bogomolov-sheaf
criterion of \cite[Theorems~4.14--4.15]{CampanaSurvey}.

\begin{proposition}[Finite stability]\label{in:finite-stability}
Suppose a finite orbifold \'{e}tale morphism has smooth connected
projective coarse source and target, and that both orbifold
boundaries are empty or disjoint unions of smooth order-two
divisors.  If the target is special in the line-bundle convention,
so is the source.  This applies both to ordinary finite \'{e}tale
pullbacks of the pair and to finite manifold covers of
\(\mathcal X\) in the orbifold sense.
\end{proposition}

\begin{proof}
Campana's \cite[Proposition~9.11]{CampanaOrbifolds}, with orbifold
\'{e}taleness as in Definition~2.23 and Proposition~2.24 there,
preserves specialness for smooth connected geometric orbifolds in
Fujiki's class \(\mathcal C\) of the same dimension, under a
morphism \'{e}tale in codimension one in the orbifold sense.  In our
application these orbifolds are projective, integral, and smooth
with simple normal crossing support.  The local condition at a
source prime is
\[
 e_E\,m_{\mathrm{source}}(E)
       =m_{\mathrm{target}}(v(E)).
\]
Over a target order-two divisor the possibilities are
\((e_E,m_{\mathrm{source}}(E))=(1,2)\) or \((2,1)\);
off the boundary they are \((1,1)\).  Thus every hypothesis of that
proposition holds.  The preceding identification of symmetric
differentials identifies its specialness criterion with ours.
\end{proof}

\section{Open fibre images, inertia, and the Albanese kernel}
\label{in:inertia}

All group images below include a choice of path to the fixed base
point.  Changing that path conjugates the image.  We use actual
images of groups of open fibres, rather than abstract groups of
compact fibres.

\begin{lemma}[Complement groups and dominant maps]\label{in:topology}
If \(U\) is a dense Zariski open in a smooth connected complex
variety \(M\), then \(\pi_1(U)\to\pi_1(M)\) is surjective and its
kernel is normally generated by meridians of the divisorial
components of \(M\setminus U\).  A dominant morphism between smooth
connected algebraic varieties has finite-index image on fundamental
groups.  The same conclusion holds for a rational map on any dense
open where it is defined.
\end{lemma}

\begin{proof}
Perturb loops and null-homotopy discs transversely to smooth strata
of the deleted set.  Loops miss it; discs miss the strata of complex
codimension at least two and meet divisors in finitely many
transverse points.  Removing small discs about these points proves
normal generation.  This argument also applies in analytic balls.
For the second assertion, take the closure of a closed point of the
generic fibre of the dominant map.  It is generically finite over
the target.  After shrinking it and the target, separability and
generic finiteness give a finite \'{e}tale cover of a target open.
The resulting finite-index subgroup maps to a finite-index subgroup
of the full target group, by the first assertion, and its image is
contained in the image under consideration.
\end{proof}

In particular, local root charts and van Kampen give
\begin{equation}\label{in:orbifold-presentation}
 \pi_1^{\mathrm{orb}}(\mathcal X)
 =\pi_1(X\setminus D)
     /\langle\!\langle\gamma_j^2\rangle\!\rangle,
\end{equation}
where one takes the meridians of all components.  Equivalently,
ordinary covers of the complement extend in root charts precisely
when the squared meridians act trivially.  This description also
identifies subgroups with connected orbifold covers.

\begin{lemma}[Open fibre images]\label{in:open-images}
Let \(f:X\dashrightarrow Z\) be a rational fibration onto a smooth
projective model.  Resolve its graph, remove the inverse image of
\(D\), and take suitable dense opens of source and base.  The image
\(B_f\) of the general open fibre group in \(G\) is well defined up
to conjugacy, is normal in \(G\), and is unchanged by further
shrinking or resolution.  There is a surjection
\begin{equation}\label{in:base-monodromy}
 \psi_f:\pi_1(Z^\circ)\longrightarrow G/B_f
\end{equation}
compatible with the map from the open total space.
The same assertions hold for any fixed homomorphic image of the
open source group.
\end{lemma}

\begin{proof}
Resolve the graph and its deleted locus.  By generic smoothness,
after shrinking the base the compact family is smooth proper and
its deleted divisor is a relative simple normal crossing divisor;
include the exceptional loci among the deletions.  The family of
pairs is differentiably locally trivial.  One can see this directly
by lifting base vector fields tangent to all boundary components
in relative crossing coordinates, patching with a partition of
unity, and using properness to integrate the lifts.  The complement
therefore is a locally trivial fibration with connected fibres.
Its homotopy sequence ends in
\[
 \pi_1(F^\circ)\longrightarrow\pi_1(X^\circ)
       \longrightarrow\pi_1(Z^\circ)\longrightarrow1.
\]
The open total space maps onto \(G\), since it contains an open
identified with a dense open of \(X\setminus D\).  Exactness gives
normality and \eqref{in:base-monodromy}.  To compare models, pass
to common dense opens of their smooth general fibres.
Lemma~\ref{in:topology} says that deleting further analytic subsets
does not change their images in the fixed target group.  Transport
in the family handles the base-point choices.
\end{proof}

\begin{lemma}[Actual inertia]\label{in:actual-inertia}
For the quotient \(R=G/B_f\) and the monodromy
\eqref{in:base-monodromy}, the image of a meridian about every prime
\(J\) on any smooth projective model of \(Z\) has finite order
\(m_J\).  Set \(m_J=1\) on primes meeting the chosen open.  The divisor
\begin{equation}\label{in:boundary}
 \Delta_Z=\sum_J(1-m_J^{-1})J
\end{equation}
has finite support on each model, and orders agree on strict
transforms.  Smooth projective models whose divisorial deleted
locus has simple normal crossings can be chosen above any given
model, and over a fixed torus whenever the map to it is already
holomorphic.
\end{lemma}

\begin{proof}
Transfer the homomorphism through common base opens and resolve the
source map to the chosen model.  A source prime dominating \(J\)
has positive pullback multiplicity \(e\).  A transverse punctured
disc at its general point lies in the compatibility open and winds
\(e\) times about \(J\).  Its loop has finite image in \(G\): the
resolved source maps holomorphically to \(X\), and a small ball in
\(X\) meets at most one smooth boundary branch.  The local
complement group has image of order at most two in \(G\).  Thus
the \(2e\)-th power of the meridian of \(J\) dies in \(R\).
This reasoning includes exceptional primes on either model.
Only deleted divisors can have nontrivial inertia, giving finite
support.  Strict transforms have the same transverse loops.
Log resolution of the complement, together with the graph of the
map to the fixed torus when needed, gives the final assertion.
\end{proof}

\begin{lemma}[Weighted local divisibility]\label{in:local-divisibility}
In the setting of Lemma~\ref{in:actual-inertia}, suppose
\(p:Z\to W\) is smooth proper with connected fibres, both spaces
are smooth, and \(\Delta_Z=p^*\Delta_W\).  Let a resolution of
\(f\) give \(u:\widehat X\to W\).  If the strict transform \(E\)
of a prime of \(X\) dominates a boundary prime \(J\) of \(W\), then
\begin{equation}\label{in:divisibility}
 m_J\mid r_E\operatorname{ord}_E(u^*J).
\end{equation}
The image of \(E\) in \(Z\) need not dominate \(p^{-1}J\).
\end{lemma}

\begin{proof}
Work over a general point of \(J\), where it is smooth and no
other boundary component of \(W\) occurs.  In a small ball about
the relevant image point in \(Z\), \(p^{-1}J\) is smooth and is
the only divisor with potentially nontrivial meridian image.
Adding back all other deleted divisors kills normally generating
meridians which already map trivially; codimension-two deletions
do not affect this calculation.  Thus the local monodromy factors
through the complement of this single smooth hypersurface.
Its generator has order \(m_J\) in \(R\).  A transverse disc to
\(E\) winds \(\operatorname{ord}_E(u^*J)\) times about it.  Since
\(E\) is nonexceptional over \(X\), its source meridian has
\(r_E\)-th power trivial in \(G\).  Compatibility proves
\eqref{in:divisibility}.
\end{proof}

\begin{lemma}[Vanishing in root coordinates]\label{in:weighted-vanishing}
Let \(u:X\to W\) be a dominant holomorphic map with
\(b=\dim W>0\).  If \(E\) dominates a smooth divisor \(J\) and
\(e=\operatorname{ord}_E(u^*J)\), the pullback of a local volume
form on \(W\) vanishes to order at least \(r_Ee-1\) at the chart
prime over \(E\).  Consequently it absorbs the pole of coefficient
\(1-1/m_J\) whenever \(m_J\le r_Ee\).  If \(u(E)\) has
codimension at least two, the pulled-back volume form vanishes at
least once at that chart prime.
\end{lemma}

\begin{proof}
The normal equation of \(J\) pulls back to a unit times
\(w_1^{r_Ee}\).  Differentiation and wedging give the first bound,
and
\(r_Ee-1\ge r_Ee(1-1/m_J)\) gives the cancellation in each
divisible pluridegree.  In the second case the tangential rank on
the chart divisor is at most \(b-2\).  Adding the transverse
direction still leaves rank at most \(b-1\) along that divisor,
so every coefficient of the pulled-back top form vanishes there.
\end{proof}

\subsection{Removing the finite-cover cases}
\label{in:finite-cover-cases}
Initially \(D\) is smooth and connected.  By
Lemma~\ref{in:root-bound}, \(X\) is special.  The manifold theorem
in Assumption~\ref{in:input-manifold} therefore makes
\(\pi_1(X)\) virtually abelian.  If the meridian \(\gamma\) is
trivial in \(G\), then \eqref{in:orbifold-presentation} and
Lemma~\ref{in:topology} identify \(G\) with \(\pi_1(X)\), settling
this case.

If \(\gamma\) survives in a finite quotient of \(G\), take the
associated normal cover of \(X\setminus D\).  Its extension across
\(D\) has the smooth local form \(z_1=w_1^2\): normality makes
every inertia orbit have length two, and connectedness of \(D\)
makes all its meridians conjugate.  The resulting connected cover
\(Y\to X\) is finite and smooth projective, by normalization and
Riemann existence.  It is orbifold \'{e}tale with empty source
boundary.  The subgroup identification for
\eqref{in:orbifold-presentation} gives
\(\pi_1(Y)=\ker(G\to F)\), where \(F\) is the chosen finite image:
the meridians filled in upstairs are exactly the lifted squares,
with all their sheet transports.  Proposition~\ref{in:finite-stability}
and Assumption~\ref{in:input-manifold} now settle this case as well.

For the remainder assume \(\gamma\ne1\) and that it dies in every
finite quotient.  Put
\[
 K=\ker\bigl(G\longrightarrow\pi_1(X)\bigr).
\]
It is the normal closure of \(\gamma\).  Every finite-index subgroup
of \(G\) contains \(K\), since its finite-index normal core does.
Choose a connected finite \'{e}tale cover of \(X\) whose ordinary
fundamental group is free abelian; such a subgroup exists in any
finitely generated virtually abelian group.  Pull up the root
datum.  The new group is the preimage of this subgroup in \(G\),
and its coarse kernel is the same \(K\).  This follows either from
orbifold covering theory or by transporting all meridian-square
relations through the ordinary cover.  The pair stays special by
Proposition~\ref{in:finite-stability}.  Relabel these data by
\(X,D,\mathcal X,G\).  The support of \(D\) is now allowed to be
disconnected but remains a disjoint union of smooth divisors.

Assumption~\ref{in:input-special} gives a surjective Albanese map
\(\alpha:X\to A\) with connected fibres.  The abelian variety
\(A=\operatorname{Alb}(X)\) has lattice
\(\Lambda=\pi_1(A)=H_1(X,\mathbb Z)/\mathrm{torsion}\).
Our ordinary group is free abelian, so the Albanese map identifies
\(\pi_1(X)\) with \(\Lambda\).  Hence
\begin{equation}\label{in:lattice-extension}
 1\longrightarrow K\longrightarrow G
       \longrightarrow\Lambda\longrightarrow1.
\end{equation}

\begin{proposition}[The Albanese kernel]\label{in:kernel}
For a general point \(a\in A\), put \(P_a=\alpha^{-1}(a)\).
Then
\begin{equation}\label{in:kernel-image}
 K=\operatorname{im}\bigl(\pi_1(P_a\setminus D)\longrightarrow G\bigr).
\end{equation}
In particular \(K\) is finitely generated.
\end{proposition}

\begin{proof}
Let \(B\) be the open fibre image in
Lemma~\ref{in:open-images}, applied to \(\alpha\).  Then \(B\subset K\),
and the induced monodromy \(\pi_1(A^\circ)\twoheadrightarrow G/B\)
has the finite actual inertia orders of
Lemma~\ref{in:actual-inertia}.  Denote its boundary by \(\Delta_A\).
If \(A\) is a point, \eqref{in:kernel-image} is immediate, so assume
otherwise.

We recall the needed effective-divisor fact on abelian varieties.
Any nonzero effective rational divisor \(\Theta\) on \(A\) is
\(p^*\Theta_U\) for a quotient \(p:A\to U\) with connected fibres
and an ample effective rational divisor \(\Theta_U\), with
\(\dim U>0\).  To see this, clear denominators and let \(L\) be
the corresponding effective line bundle.  The connected kernel
of \(x\mapsto t_x^*L\otimes L^{-1}\) acts projectively on
\(H^0(A,L)\).  Its action is trivial, since a morphism from a
connected complete group to the affine group \(\operatorname{PGL}\)
is constant.  Thus it fixes the divisor, which descends to the
quotient by this connected kernel.  The descended effective line
bundle is nef (translate its divisor to test intersection with
any curve), and its translation kernel is finite.  The line-bundle
criterion on abelian varieties makes it ample: the associated
semipositive Hermitian form has radical the connected translation
kernel, so here is positive definite.  The quotient has positive
dimension because the divisor is nonzero.  Prime components and
their rational coefficients pull back unchanged under the smooth
quotient.

Suppose now \(\Delta_A\ne0\), and write
\(\Delta_A=p^*\Delta_U\) in this way.  Put \(b=\dim U\) and
\(u=p\alpha\).  Let \(\mathcal L_u\subset\Omega^b_{\mathcal X}\)
be the saturated line of pulled-back top forms.  For every prime
\(E\) dominating a boundary prime \(J\) of \(U\),
Lemma~\ref{in:local-divisibility} gives
\(m_J\mid r_E\operatorname{ord}_E(u^*J)\).  Thus
Lemma~\ref{in:weighted-vanishing} cancels the full allowed pole
along \(E\).  For a prime mapping into codimension at least two
in \(U\), that lemma gives positive vanishing.  There are only
finitely many primes in \(u^*\Delta_U\); choose a single rational
\(0<\epsilon\le1\) small enough that this positive vanishing
cancels the coefficient of \(u^*(\epsilon\Delta_U)\) at each of
these remaining primes, in root coordinates.

For every sufficiently divisible \(j\), pullback consequently gives
\[
 H^0\bigl(U,j(K_U+\epsilon\Delta_U)\bigr)
       \longrightarrow H^0(\mathcal X,\mathcal L_u^{\otimes j}).
\]
This is checked at chart primes and extends across codimension
two.  Dominance preserves section ratios.  Since \(K_U\) is
trivial and \(\epsilon\Delta_U\) is ample, the right-hand line
has Iitaka dimension at least \(b\), hence exactly \(b\) by
Lemma~\ref{in:root-bound}.  This contradicts specialness.

It follows that \(\Delta_A=0\).  All deleted divisorial meridians
are therefore killed, and Lemma~\ref{in:topology} factors the base
monodromy through \(\pi_1(A)\).  Its induced map to \(G/B\) is
inverse to the natural map \(G/B\to\Lambda\): both composites
are determined on the surjecting open groups.  Hence \(B=K\).
The smooth quasi-projective variety \(P_a\setminus D\) has finite
CW type, so its fundamental group, and its image \(K\), are
finitely generated.
\end{proof}

\begin{lemma}[Finite kernels over a lattice]\label{in:finite-by-lattice}
If \(1\to F\to E\to\mathbb Z^r\to1\) is exact with \(F\) finite,
then \(E\) has a finite-index free abelian subgroup disjoint from
\(F\).  If \(K\) is a finitely generated normal subgroup of a group
\(G\) with lattice quotient, every finite-index \(K_0\le K\)
contains \(K\cap G_0\) for some finite-index \(G_0\le G\).
\end{lemma}

\begin{proof}
Centralize \(F\) by passing to a finite-index subgroup \(E_1\).
Its commutator subgroup lies in the finite central group
\(F\cap E_1\).  If \(x_1,\ldots,x_t\) generate \(E_1\), and
\(e\) annihilates \(F\), then
\([x_i^e,x_j^e]=[x_i,x_j]^{e^2}=1\).  The subgroup generated by
these powers is abelian and has finite index: its image has finite
index in the abelianization, whose kernel here is finite.  A
torsion-free subgroup of finite index in this finitely generated
abelian group proves the first claim.

For the second, a finitely generated group has only finitely many
subgroups of each specified index.  Intersect all subgroups of
\(K\) having index \([K:K_0]\).  The result \(K_1\subset K_0\)
is characteristic and finite-index in \(K\), hence normal in
\(G\).  Apply the first claim to \(G/K_1\), whose finite kernel
is \(K/K_1\), and pull back a subgroup disjoint from this kernel.
\end{proof}

\begin{proposition}[The remaining group and its covers]
\label{in:group-reduction}
In the remaining case of \eqref{in:lattice-extension}, the group
\(K\) is infinite, finitely generated, and has no proper
finite-index subgroup.  All finite-dimensional complex linear
representations of \(K\), or of any quotient of \(K\), are trivial.
Every finite-index subgroup of \(G\) is the preimage of a
finite-index sublattice of \(\Lambda\), and is geometrically
realized by pulling the pair back along the corresponding
isogeny of \(A\).  These pullbacks are connected and special;
their kernels are the same \(K\), and the open fibre images in
\eqref{in:kernel-image} identify with the original ones.
\end{proposition}

\begin{proof}
We already know that every finite-index subgroup of \(G\)
contains \(K\).  Proposition~\ref{in:kernel} and
Lemma~\ref{in:finite-by-lattice} therefore rule out every proper
finite-index subgroup of \(K\).  The group \(K\) is nontrivial
because it contains the surviving meridian; it must be infinite,
since a nontrivial finite group has a proper finite-index trivial
subgroup.  A complex linear image is finitely generated and
residually finite by Malcev's theorem; see \cite{Nica}.  Any
nonidentity element in such an image would survive in a finite
quotient of \(K\), which is impossible.  The same reasoning
applies to its quotients.

If \(G_0\le G\) has finite index, then \(K\subset G_0\), so
\(G_0\) is exactly the preimage of its lattice image
\(\Lambda_0\).  The inclusion \(\Lambda_0\subset\Lambda\)
defines an isogeny \(A_0\to A\).  The pullback
\(X_0=X\times_A A_0\) is connected: on ordinary fundamental
groups the map \(\pi_1(X)\to\Lambda\) is an isomorphism, so
the monodromy on the isogeny's fibre is transitive.  Its pulled-back
root orbifold has group \(G_0\), by the complement presentation,
and is special by Proposition~\ref{in:finite-stability}.  A fibre
over a specified point of \(A_0\) is isomorphic to the original
fibre over its image, including the boundary.  These identifications
give the asserted equality of open images inside \(K\); images
of resolved subvarieties are compared on common dense opens in
the same way.
\end{proof}
